\section{Introduction}
\label{sec:intro}

A bank that declines a loan, or a clinic that flags a patient, is often asked for the reasons. For a tabular model the usual answer is the few attributes with the largest Shapley attributions, obtained by running a sampling-based estimator such as KernelSHAP \citep{lundberg2017unified} once on that instance, at the library's default budget, and reporting the top of the list. The estimator is random and its budget is limited, so a second run may report different reasons. Two attributions with nearly equal magnitude may swap, and a systematic error in the estimator may place the wrong attribute in the set in every run. Explanations are checked rarely, and when they are, it is usually by the mean error over many instances.

We ask when a single run is reliable evidence for the explanation of \emph{that} prediction, and whether a deployer can tell from the run itself. We keep two things apart. \emph{Repeatability} asks whether a rerun would report the same attributes; \emph{accuracy} asks whether the run reports what the exact Shapley value would. The first can be measured without the exact value, but the second is what a reason is for, and an estimator can be repeatable without being accurate.

Work on Shapley estimation bounds errors of the full attribution vector \citep{chen2023algorithms,musco2024leverage}, and a few methods give guarantees for identifying the top-$k$ attributes and stop sampling when these hold \citep{covert2021improving,kariyappa2024shapk,goldwasser2024significance,kolpaczki2025antithetic}. Studies of explanation stability and disagreement concern the variation of the explained quantity itself, such as the choice of background \citep{yuan2022background} or of another explainer \citep{krishna2022disagreement}. What is missing for tabular models is a decision-level, per-instance measurement of how often one run of the \emph{shipped} estimator reports the right top-$k$ set, how this depends on the model, the dimension and the budget, and whether the in-run certificates that exist can be trusted when used as intended.

We provide this measurement. Our approach has three parts. First, we measure reliability per instance and at the decision level: for each instance, the probability $r_k$ that a run returns exactly the top-$k$ set of the exact interventional Shapley value, estimated from repeated runs, and the fraction of instances for which this probability is at least $1-\delta$; the reference is computed exactly by enumeration. Second, we read these numbers against explicit scales: chance, the gap between the $k$-th and $(k{+}1)$-th attribution relative to the estimator's standard error, and the variation that practitioners already accept from resampling the background and from retraining. Third, we audit certificates from the literature for calibration on fresh reruns, under the stopping rule they prescribe.

On two public tasks, three model classes and three seeds, we find:
\begin{itemize}
\item \textbf{The mean hides a tail, and a default fails repeatedly.} On ACS Employment ($d=16$) the shipped KernelSHAP returns the exact top-3 set in 97\% of runs on average, but 7--10\% of instances are unreliable; these are near ties, and the top-3 is correct up to a 10\% tolerance in essentially every run. On German Credit ($d=20$) the same default is correct in only 83--88\% of runs. The error is repeatable (rerun stability 0.95--0.97), it is not cured by sixteen times the budget, and it comes from the default's $\ell_1$ feature selection; switching selection off at the same budget gives 97--99\%.
\item \textbf{Certificates are calibrated but conservative, and dear.} Four of five certification rules deliver at least their nominal level on fresh reruns; the fifth, which tests the boundary pair on the run's own estimate, over-claims for the larger $\delta$ and $k$ in three of six cells. The calibrated rules abstain on 22--76\% of the instances an oracle would certify at four times the default budget, and cost at least about forty times what a fixed corrected estimator needs for the same exact rate (the fixed estimator is at the floor of our budget grid, and makes no per-run statement).
\end{itemize}
These findings give the verdict its form: one run suffices for a corrected estimator at a small budget on these models, not for the shipped default at $d=20$, and which instances of the remaining few percent are affected cannot be read from the run without a certificate.

Our contribution is the audit and its evidence, not a new estimator or certificate. We release the code, the exact reference values and all result files with this submission.
